\documentclass{plantillaPPT}
\usepackage{esint}
\titulo{13}{Teoremas Importantes y Fin de Cálculo III}
\begin{document}
\primeraDiapo

\begin{frame}{Teorema de Stokes}
\framesubtitle{Teorema}
    \begin{teorema}
        Sea $S$ una superficie suave (o suave por tramos) con borde y orientable, cuyo borde $C=\partial S$ es una curva simple (no se corta a si misma), cerrada y orientada positivamente respecto a $S$.
        Si $\vec{F}$ es $C^1$ en un abierto $A$ que contiene a $S\cup C$, entonces
        \[
        \oint\limits_C \vec{F}\cdot d\vec{r} = \iint\limits_S \nabla\times\vec{F}~d\vec{S}
        \]
    \end{teorema}
    
\end{frame}

\begin{frame}{Problema 1}
    \framesubtitle{Práctica 14}
    1. Verifique el \textbf{Teorema de Stokes} con el campo $\vec{F}:\mathbb{R}^3\to\mathbb{R}^3$ dado por $\vec{F}(x,y,z) = (-y,x,1)$ y sea $S$ el paraboloide $z=x^2+y^2, 0\leq z\leq1$, con sus vectores normales apuntando hacia afuera de él. Evalúe
    \[
    \iint\limits_S \nabla\times\vec{F}~d\vec{S}
    \]
\end{frame}
\begin{frame}{Teorema de la Divergencia (de Gauss)}
    \framesubtitle{Teorema}
    \begin{teorema}
        Si $\Omega$ es una región del espacio 3-D acotada por la \textcolor{red}{superficie cerrada} $S=\partial\Omega$, suave (o suave por tramos) y orientada exteriormente y $\vec{F}= (P,Q,R)$ es un campo vectorial de clase $C^1$ en un abierto que contiene a $\Omega\cup S$, entonces
        \[
        \oiint\limits_S \vec{F}\cdot d\vec{S} = \iiint\limits_\Omega \nabla\cdot\vec{F}~dV
        \]
        donde {\color{blue} $\nabla\cdot\vec{F} = \frac{\partial P}{\partial x}+\frac{\partial Q}{\partial y}+\frac{\partial R}{\partial z}$}.
    \end{teorema}
    
\end{frame}
\begin{frame}{Problema 2}
    \framesubtitle{Práctica 13}
    2. Verifique el \textbf{Teorema de Gauss o de la Divergencia}
    \[
    \oiint\limits_S \vec{F}\cdot d\vec{S} = \iiint\limits_D \nabla\cdot\vec{F}~dV,
    \]
donde $\vec{F}(x,y,z) = x^2\hat{i}+xz\hat{j}+3z\hat{k}$ y el borde de la superficie $S$ de la esfera sólida $x^2+y^2+z^2\leq4$
\end{frame}


\end{document}